% Chapter 6 section 6.1: overview.
% Written 2026-09-26 (manifest §9b row 6.1), drafted in the chat and approved before this write.
% Shape follows §5.1 as rewritten 2026-09-25: no verbatim RQ block and no mapping table (§15 O-20 c),
%   no float, no ranking clause (review E4), ~250 words.
% Carries: the road map; the recap of which configuration "the multi-index retriever" names -- the
%   blocked system, as §5.6 selects it, with no numbers repeated; FC-P introduced for the chapter;
%   faithfulness stated once as attribution rather than quality (manifest §1c); and the reading rule --
%   generation differences read as reported, §5.1's 0.02 belonging to retrieval scores alone.
% Not here: the RQ text, k, the generator sizes and the prompt (all §6.2); any oracle condition
%   (not reported -- manifest §9b, 2026-09-26).
% RQ3 is named as a sanity check on RQ2, her own term for it (susanne-review-response.md F8).
% The verdicts section is described rather than \ref'd: §6.4 has no label until it is written.

\section{Overview}\label{sec:gen-overview}

This chapter answers the second research question, on generation, and the third, a sanity
check on the second: whether its answers survive a change of generator.
Section~\ref{sec:gen-method} states how the answers were produced and how they were
scored. Section~\ref{sec:gen-results} reports three experiments: how answers generated
from the multi-index retriever's context compare with answers generated from no context and
from the context of text-only retrieval; whether answer quality tracks retrieval quality
across the fifteen retrieval configurations of
Section~\ref{sec:met-retrieval-configs}; and whether both results hold as the generator
grows. The chapter closes by ruling on each hypothesis.

The retriever this chapter gives the generator is the one Section~\ref{sec:ret-verdicts}
selected, the blocked multi-index system, and the retrieval it is set against is the
strongest text-only system, dense retrieval on its own. Both were chosen there by reading the
retrieval results, and this chapter takes them as given.

Two measures are reported, and they answer different questions. Factual-correctness precision
(FC-P) is the quality measure: the share of an answer's claims that the reference supports.
Faithfulness measures attribution rather than quality --- the share of an answer's claims that
the retrieved context supplied --- so a low value on a \emph{correct} answer means the
generator already knew the answer, not that the answer is poor. Faithfulness needs a context
to be computed against, and generation without retrieval therefore has no value for it.

Retrieval quality enters as the primary measure of Chapter~\ref{ch:retrieval},
nDCG$'$@10 under the labels of the relevance judge of
Section~\ref{sec:met-judge-validation}. Differences between generation scores are read as the
values are reported, with no statistical test and no equivalent limit: the 0.02 of
Section~\ref{sec:ret-overview} is a property of the relevance judge and governs retrieval
scores alone.
